\begin{helper}[Normal witness for the rotated trunk]
Let $N$ be balanced normal globally maximum data for
$(G,k,e_0,\varphi,T)$, put $p=N.height$, fix $1\leq j<p$, and put
$f_i=\noderef{NormalCycleEdge}(N,i)$.  Let distinct $A,B$ be the parity
choices
\[
 A=\begin{cases}N.beta,&j\text{ odd},\\N.alpha,&j\text{ even},\end{cases}
 \qquad
 B=\begin{cases}N.alpha,&j\text{ odd},\\N.beta,&j\text{ even}.
 \end{cases}
\]
Suppose that $\psi(f_i)=\varphi(f_{i+1})$ for $0\leq i<j$ and that
$\psi$ fixes every identity outside $\{f_0,\ldots,f_j\}$.

Let $T'$ be a tree sequence based at $f_j$ which is globally maximum for
$\psi$ and has $T'.length=T.length$.  Suppose $A$ is missing at $y_j$,
$B$ is missing at $y_{(j+1)\bmod p}$, and the following exact trunk
formulas hold:
\[
 \begin{aligned}
 T'.vertex(r)&=T.vertex((j+r)\bmod p) &&(r<p),\\
 T'.edge(r)&=f_{(j+r)\bmod p} &&(r+1<p),
 \end{aligned}                                                   \tag{1}
\]
where each right-hand index is represented by the finite witness required
in the Lean statement.

Then there are normal globally maximum data $N'$ for
$(G,k,f_j,\psi,T')$ such that
\[
 N'.height=p,qquad N'.alpha=A,qquad N'.beta=B,
\]
and, for
$pred=(j+p-1)\bmod p\in\operatorname{Fin}(p)$,
\[
 N'.closingEdge=f_{pred}.
\]
\end{helper}
\begin{proof}
The cycle theorem \noderef{ScheideTrunkRotationCycle} gives $p\geq3$,
the exact cyclic endpoint pairs, pairwise identity distinctness, and the
old alternating colors.  Define $N'$ to have the supplied global field,
height $p$, colors $A,B$, and closing identity
\[
 q'=f_{(j+p-1)\bmod p}.                                      \tag{2}
\]
The lower height bound follows from $p\geq3$.  The old height upper bound
$p\leq T.length+1$ and the supplied equality
$T'.length=T.length$ give the new upper bound.

Formula (1) at $r=0$ identifies $T'.vertex(0)$ with $y_j$, because
$j<p$ makes $(j+0)\bmod p=j$.  At $r=1$, it identifies
$T'.vertex(1)$ with $y_{(j+1)\bmod p}$.  The two assumed missing-color
facts are therefore exactly the alpha- and beta-missing fields required
of $N'$.  The closing-cycle endpoint equation and (1) give
\[
 G.ends(q')=\{T'.vertex(p-1),T'.vertex(0)\};                   \tag{3}
\]
indeed the cyclic edge immediately preceding $f_j$ joins the last vertex
of the rotated block to its first.  Take the new last index to be $p-1$;
it lies below $T'.length+1$, has successor $p$, and (3) proves the
closing-end field.  This includes the wraparound $j=p-1$, where (2) is
$f_{p-2}$.

We now construct the stored closing walk explicitly.  In the two-color
incidence graph for $\psi,A,B$, take
\[
 \begin{split}
 R={}&(\operatorname{inl}(T'.vertex(1)),
       \operatorname{inr}(T'.edge(1)),
       \operatorname{inl}(T'.vertex(2)),\ldots,\\
    &\operatorname{inr}(T'.edge(p-2)),
       \operatorname{inl}(T'.vertex(p-1)),
       \operatorname{inr}(q'),
       \operatorname{inl}(T'.vertex(0))).                     \tag{4}
 \end{split}
\]
For $p=3$, the middle range $1\leq r\leq p-2$ has exactly one edge; the
notation in (4) therefore displays that edge once, not twice.

By the shift and fixed equations, the old alternation from
\noderef{ScheideTrunkRotationCycle}, and the parity definitions of $A,B$,
every identity displayed in (4) is colored $A$ or $B$ under $\psi$.
Formula (1), equation (3), and the cyclic endpoint equations show that
each is incident with its two neighboring displayed vertices.  Thus every
successive pair is adjacent by \noderef{TwoColorIncidenceGraph}, so (4)
is a walk from $T'.vertex(1)$ to $T'.vertex(0)$.

The $p$ displayed vertex nodes are pairwise distinct because (1) is a
cyclic permutation of the first $p$ old vertices and
\noderef{TreeSequence} has an injective vertex map.  The $p-1$ displayed
edge nodes are precisely all cycle identities except $f_j$, and are
pairwise distinct by the cycle theorem.  Sum tags separate vertex nodes
from edge nodes.  Hence $R$ is a simple path and has $2(p-1)$ incidence
steps.

The same enumeration proves both support equivalences required by
\noderef{NormalGlobalMaximumTriple}.  Its vertex-node support is exactly
\[
 \{\operatorname{inl}(T'.vertex(i)):0\leq i<p\}.
\]
Its edge-node support is exactly the closing identity $q'$ together with
\[
 \{\operatorname{inr}(T'.edge(r)):1\leq r,\ r+2\leq p\}.
\]
The forward directions follow by inspecting (4).  Conversely, every
index in the displayed bounds occurs in (4), and pairwise distinctness
prevents an omitted exact parallel identity from being identified merely
through equal endpoints.  Finally, the cyclic endpoint equation and (1)
give, for each $1\leq r\leq p-2$,
\[
 G.ends(T'.edge(r))=
 \{T'.vertex(r),T'.vertex(r+1)\},
\]
which is the normal trunk-edge endpoint field.  We have now constructed
every field of $N'$.  Its height and colors were chosen as asserted, and
(2) is exactly the finite predecessor witness in the conclusion.
\end{proof}
